\thispagestyle{firstpage}
\begin{flushleft}
{\sffamily\bfseries\fontsize{17}{20.5}\selectfont \reportassignment:\\ Instruction Tuning and Direct Preference Optimization\\ on a Small Language Model\par}
\vspace{0.8em}
{\sffamily\bfseries\normalsize \reportauthor\par}
{\sffamily\small \reportprogram, Roll No.\ \reportroll, \reportinstitution\par}
{\sffamily\small \reportcourse; instructor: \reportinstructor\par}
\end{flushleft}
\vspace{0.3em}

\noindent{\sffamily\bfseries\small ABSTRACT}\par\vspace{2pt}
{\small\noindent This report follows a small post-training pipeline on Qwen2.5-0.5B: supervised fine-tuning (SFT) on 30,000 UltraChat examples, then two Direct Preference Optimization (DPO) runs written from scratch, one on 15,000 HH-RLHF pairs and one on 15,000 UltraFeedback pairs. A third experiment trains a LoRA persona adapter with TRL. I compare each DPO model with its SFT starting point on 300 fixed prompts, using a pairwise language-model judge. The HH-RLHF model beats SFT on its own safety prompts (win rate 0.596, $p=0.004$). The UltraFeedback model loses on AlpacaEval (0.392) because its answers are longer and repeat themselves more, and the loss turns into a tie (0.498) when I discourage repetition. The persona adapter either shows no persona or gains style at the cost of usefulness and safety. One judge, small evaluation sets, and a single training seed limit how far these conclusions reach.\par}
\vspace{0.4em}
\noindent{\sffamily\bfseries\small Additional Key Words and Phrases:} {\small preference optimization, DPO, supervised fine-tuning, LoRA, LLM-as-a-judge, small language models\par}
\vspace{0.2em}
